\documentclass[12pt,a4paper]{article}
\usepackage[utf8]{inputenc}
\usepackage[ngerman]{babel}
\usepackage[T1]{fontenc}
\usepackage{geometry}
\geometry{margin=2.5cm}
\usepackage{amsmath,amsfonts,amssymb}
\usepackage{graphicx}
\usepackage{hyperref}
\hypersetup{
	colorlinks=true,
	linkcolor=blue,
	citecolor=red,
	urlcolor=blue,
}
\usepackage{enumitem}
\usepackage{setspace}
\onehalfspacing
\usepackage{csquotes}

\title{Frauenrechte in Deutschland und Europa:\\Stand, Herausforderungen und Handlungsempfehlungen}
\author{}
\date{\today}

\begin{document}
	
	\maketitle
	
	\begin{abstract}
		Die gleichberechtigte Teilhabe von Frauen in allen gesellschaftlichen Bereichen ist ein fundamentales Menschenrecht. Trotz erheblicher Fortschritte in den letzten Jahrzehnten besteht in Deutschland und Europa nach wie vor erheblicher Handlungsbedarf. Das vorliegende Dokument analysiert den aktuellen Stand der Frauenrechte in Deutschland im europäischen Vergleich, identifiziert zentrale Herausforderungen — insbesondere im digitalen Raum und im Kontext Künstlicher Intelligenz — und formuliert darauf aufbauend politische Handlungsempfehlungen. Ein besonderer Fokus liegt auf der Rolle von Männern als Verbündete sowie auf den Potenzialen geschlechtergerechter Bildung als präventiver und transformativer Strategie.
	\end{abstract}
	
	\newpage
	
	\section{Einleitung}
	
	Die gleichberechtigte Teilhabe von Frauen in allen gesellschaftlichen Bereichen ist ein fundamentales Menschenrecht und ein zentraler Indikator für den demokratischen Entwicklungsstand einer Gesellschaft. Trotz erheblicher Fortschritte in den letzten Jahrzehnten besteht in Deutschland und Europa nach wie vor erheblicher Handlungsbedarf. Das vorliegende Dokument analysiert den aktuellen Stand der Frauenrechte in Deutschland im europäischen Vergleich, identifiziert zentrale Herausforderungen — insbesondere im digitalen Raum — und formuliert darauf aufbauend politische Handlungsempfehlungen.
	
	\section{Aktuelle Lage der Frauenrechte in Deutschland}
	
	\subsection{Fortschritte und verbleibende Defizite}
	
	Die Europäische Union hat in den letzten Jahrzehnten erhebliche Fortschritte bei der Gleichstellung der Geschlechter erzielt, unter anderem durch Rechtsvorschriften, Gender-Mainstreaming und besondere Maßnahmen zur Förderung von Frauen. Mehr Frauen sind berufstätig und erlangen höhere Bildungsabschlüsse als je zuvor. Deutschland hingegen hinkt im europäischen Vergleich in zentralen Gleichstellungsindikatoren hinterher.
	
	\subsection{Gewaltschutz und die Istanbul-Konvention}
	
	Die Istanbul-Konvention des Europarats zur Verhütung und Bekämpfung von Gewalt gegen Frauen und häusliche Gewalt trat in Deutschland 2018 in Kraft. Ein wichtiger Schritt war der Beschluss des Gewalthilfegesetzes im Februar 2025, wonach Frauen ab 2032 einen kostenfreien Rechtsanspruch auf Schutz und Beratung erhalten. Dieses Gesetz markiert einen historischen Schritt im Schutz vor geschlechtsspezifischer und häuslicher Gewalt, da es erstmals eine bundesweite gesetzliche Grundlage für ein bedarfsgerechtes Hilfesystem schafft \cite{gewalthilfegesetz}.
	
	Gleichzeitig zeigt der Alternativbericht des Bündnisses Istanbul-Konvention (BIK) von November 2025 massive Lücken beim Schutz vor geschlechtsspezifischer und häuslicher Gewalt auf \cite{bik_alternativbericht}. GREVIO — das unabhängige Expertengremium des Europarats — stuft die Umsetzung in Deutschland als unzureichend ein. Insbesondere vulnerablere Gruppen wie Frauen mit Behinderungen, Geflüchtete, wohnungs- und obdachlose Frauen sowie trans-, inter- und queere Personen haben nach wie vor kaum Zugang zum Hilfesystem.
	
	\subsection{Arbeitsmarkt und Gender Pay Gap}
	
	Der unbereinigte Gender Pay Gap in Deutschland lag 2025 unverändert bei 16 Prozent. Frauen verdienten mit 22,81 Euro pro Stunde durchschnittlich 4,24 Euro weniger als ihre männlichen Kollegen (27,05 Euro). Diese Lohnlücke blieb im Vergleich zum Vorjahr konstant \cite{gpg_2025}.
	
	Auch bei vergleichbarer Qualifikation und Tätigkeit bleibt ein bereinigter Gender Pay Gap von rund sechs Prozent bestehen. Diese verbleibende Lohndifferenz von 1,71 Euro pro Stunde lässt sich statistisch nicht eindeutig erklären und wird unter anderem auf Diskriminierung auf dem Arbeitsmarkt zurückgeführt \cite{gpg_bereinigt}.
	
	Regional fallen die Unterschiede drastisch aus: Während die Lohnlücke in Ostdeutschland lediglich fünf Prozent beträgt, liegt sie in Westdeutschland bei 17 Prozent \cite{gpg_ost_west}.
	
	\subsection{Führungspositionen}
	
	Die Unterrepräsentanz von Frauen in Führungspositionen ist eines der hartnäckigsten Gleichstellungsprobleme. Seit 2014 hat sich der Frauenanteil in Führungspositionen in Deutschland praktisch nicht verändert (+0,1 Prozentpunkte), während er im EU-Durchschnitt um 3,4 Prozentpunkte gestiegen ist. Spitzenreiter ist Schweden mit einem Frauenanteil von 44,4 Prozent.
	
	\subsection{Politische Partizipation}
	
	Im Deutschen Bundestag gehören von 630 Abgeordneten nur 204 Frauen an, was einem Anteil von 32,4 Prozent entspricht. Nach der Bundestagswahl 2025 ist der Frauenanteil damit wieder gesunken, nachdem 2017 mit 37,3 Prozent ein Höchststand erreicht worden war. Deutschland rutschte im weltweiten Ranking nationaler Parlamente von Platz 45 auf Platz 58 ab \cite{frauenanteil_bundestag}.
	
	\section{Digitale Räume als besondere Herausforderung für die Gleichstellung}
	
	\subsection{Digitale Gewalt als strukturelles Problem}
	
	Digitale Gewalt gegen Frauen ist kein Randphänomen, sondern eine strukturelle Form geschlechtsspezifischer Gewalt. Sie umfasst Cyberstalking, Cybermobbing, Doxing, Deepfakes, Sextortion und die nicht-einvernehmliche Verbreitung intimer Aufnahmen. Neue Technologien eröffnen Tätern nahezu unbegrenzte Zugriffsmöglichkeiten auf Betroffene, da Gewalt jederzeit und von überall ausgeübt werden kann.
	
	Der Digital Services Act (DSA), der seit Februar 2024 gilt, harmonisiert die Sorgfaltspflichten für Online-Plattformen hinsichtlich illegaler Inhalte. Der Schutz vor digitaler Gewalt bleibt jedoch weiterhin lückenhaft: Der DSA bietet weder Löschungsansprüche für Betroffene noch Fristen für das Tätigwerden von Diensteanbietern. Problematisch ist zudem, dass der DSA nicht definiert, was "rechtswidrige" Inhalte sind, und die Definition damit den einzelnen Nationalstaaten überlässt \cite{dsa_luecken}.
	
	\subsection{Der Silencing-Effekt}
	
	Digitale Gewalt führt zu einem massiven \textit{Silencing-Effekt}: Frauen und Mädchen ziehen sich aus Angst vor Angriffen aus sozialen Medien zurück, verzichten auf öffentliche Meinungsäußerungen oder geben politische Ämter auf. Wo Frauen sich im Netz öffentlich äußern, riskieren sie sexistische Belästigung, pornografische Pöbeleien, die Androhung von Vergewaltigungen bis hin zu Morddrohungen. Viele Frauen ziehen sich zurück und partizipieren nicht mehr am digitalen öffentlichen Diskurs, was nicht nur die Verletzung individueller Rechte bedeutet, sondern auch der Demokratie schadet \cite{dsa_luecken}.
	
	\subsection{Rechtliche Rahmenbedingungen und Schutzlücken}
	
	Deutschland ist völker- und europarechtlich verpflichtet, geschlechtsspezifische Gewalt auch im digitalen Raum wirksam zu verhindern, zu verfolgen und zu sanktionieren. Maßstäbe ergeben sich aus CEDAW, der EMRK, der Istanbul-Konvention, dem Digital Services Act (DSA) sowie der EU-Gewaltschutz-Richtlinie 2024.
	
	Die Richtlinie (EU) 2024/1385 des Europäischen Parlaments und des Rates zur Bekämpfung von Gewalt gegen Frauen und häuslicher Gewalt trat am 13. Juni 2024 in Kraft \cite{eu_gewalt_richtlinie}. Sie kriminalisiert unter anderem die nicht-einvernehmliche Weitergabe von intimem Material, Cyberstalking und Cybermobbing \cite{eu_gewalt_richtlinie}. Deutschland muss die Richtlinie bis zum 14. Juni 2027 in nationales Recht umsetzen \cite{eu_gewalt_umsetzung}.
	
	\section{Künstliche Intelligenz als Chance und Risiko für die Gleichstellung der Geschlechter}
	
	\subsection{Algorithmische Diskriminierung als strukturelles Problem}
	
	Künstliche Intelligenz (KI) gilt häufig als objektive Entscheidungsinstanz. Tatsächlich sind KI-Systeme jedoch nicht neutral, sondern spiegeln die gesellschaftlichen Vorurteile wider, die in ihren Trainingsdaten enthalten sind. Eine aktuelle Studie der Technischen Hochschule Würzburg-Schweinfurt (THWS) aus dem Juli 2025 zeigt, dass große Sprachmodelle wie ChatGPT Frauen für Gehaltsverhandlungen durchgängig niedrigere Zielbeträge empfehlen als Männern – selbst wenn alle anderen Faktoren identisch sind \cite{thws_gehaltsbias}. „Gerade bei sensiblen Themen wie Gehalt kann diese Form von verstecktem Bias reale Auswirkungen auf die Lebensrealität von Nutzerinnen haben“, warnt Prof. Dr. Ivan Yamshchikov \cite{thws_gehaltsbias_pm}.
	
	Besonders problematisch ist, dass solche Verzerrungen vor allem in realitätsnahen, interaktiven Szenarien auftreten, während Standard-Benchmarks oft keine signifikanten Unterschiede erkennen lassen \cite{thws_gehaltsbias}. Die Studie unterstreicht die Dringlichkeit ethischer Leitlinien für KI-Assistenten, wie sie etwa im EU-geförderten Projekt AIOLIA entwickelt werden \cite{thws_gehaltsbias_pm}.
	
	\subsection{Geschlechtsspezifische digitale Gewalt durch KI}
	
	Eine besonders alarmierende Entwicklung ist die Nutzung von KI zur Erstellung sexualisierter Deepfakes. Nach Angaben des Bundeskriminalamts (BKA) sind 98 Prozent aller online verfügbaren Deepfakes sexualisierte Inhalte, und 99 Prozent der Betroffenen sind weiblich \cite{bka_deepfakes}. Die Erstellung solcher Inhalte erfordert keinerlei technische Vorkenntnisse mehr – ein einziges Foto genügt, um in weniger als 25 Minuten ein täuschend echtes pornografisches Deepfake zu erstellen \cite{messer_deepfake}.
	
	Die Auswirkungen auf die Betroffenen reichen von Angstzuständen bis hin zu Depressionen und Suizid \cite{bka_deepfakes}. Besonders betroffen sind Frauen in der Öffentlichkeit, etwa Politikerinnen, Wissenschaftlerinnen und Aktivistinnen \cite{bka_deepfakes}. Der sogenannte \textit{Silencing-Effekt} – der Rückzug aus der öffentlichen digitalen Kommunikation aus Angst vor solchen Angriffen – gefährdet nicht nur die individuelle Teilhabe, sondern die demokratische Meinungsvielfalt insgesamt.
	
	Die strafrechtliche Verfolgung dieser Taten ist derzeit unzureichend: Die Herstellung nicht-einvernehmlicher sexualisierter Deepfakes ist in Deutschland überwiegend strafrechtlich nicht erfasst, und die Verfolgung scheitert häufig an mangelnder Sensibilisierung der Behörden sowie an Opferbeschuldigung (\textit{Victim blaming}) \cite{bka_deepfakes}.
	
	\subsection{Regulatorische Rahmenbedingungen und feministische Gegenstrategien}
	
	Auf europäischer Ebene wurden wichtige regulatorische Grundsteine gelegt. Mit dem \textit{EU Artificial Intelligence Act (EU AI Act)} wurde eine Grundlage für die Regulierung von KI in der EU geschaffen, die das Potenzial für eine wertebasierte, inklusive und diskriminierungsfreie Gestaltung der KI bietet \cite{eu_ai_act_frauenrat}.
	
	Der Deutsche Frauenrat fordert in einem Beschluss vom 27. Juni 2025 die Bundesregierung auf, eine inklusive, feministische KI-Strategie zu entwickeln, die unter anderem:
	\begin{itemize}
		\item verbindliche Geschlechterverträglichkeitsprüfungen für KI-Systeme einführt,
		\item durch diverse Teams eine inklusive KI-Entwicklung fördert,
		\item den Gender Data Gap (Unterrepräsentation von Frauen in Trainingsdaten) beendet und
		\item die Beteiligung von Frauenverbänden in KI-Ethikräten sicherstellt \cite{eu_ai_act_frauenrat}.
	\end{itemize}
	
	Gleichzeitig warnt der Frauenrat vor einem internationalen Rollback bei Gleichstellungsstandards, insbesondere angesichts der Anti-DEI-Dekrete der US-Regierung (DEI: Diversity, Equality/Equity, Inclusion) \cite{eu_ai_act_frauenrat}.
	
	Der Europarat beabsichtigt darüber hinaus, eine Empfehlung zur Gleichstellung und künstlichen Intelligenz anzunehmen, die die Mitgliedstaaten dabei unterstützen soll, die Gleichstellung, einschließlich der Geschlechtergleichstellung, in allen Phasen des KI-Lebenszyklus zu fördern und Diskriminierung zu bekämpfen \cite{europarat_empfehlung_ki}.
	
	\subsection{Fazit: KI als feministische Gestaltungsaufgabe}
	
	Künstliche Intelligenz ist kein neutrales Werkzeug, sondern ein gesellschaftlicher Aushandlungsraum, in dem bestehende Machtverhältnisse reproduziert oder transformiert werden können. Ohne gezielte regulierende Eingriffe droht KI, bestehende Geschlechterungleichheiten zu verfestigen und zu verstärken. Mit einer konsequenten, feministisch ausgerichteten Regulierung – allen voran durch den EU AI Act – und verbindlichen Gleichstellungsstandards bietet KI jedoch auch die Chance, eine inklusivere und gerechtere digitale Zukunft zu gestalten.
	
	\section{Männer als Verbündete und Bildungsreformen als Schlüssel zur Gleichstellung}
	
	Die bislang betrachteten Maßnahmen richten sich überwiegend an staatliche und institutionelle Akteure. Ein nachhaltiger Wandel der Geschlechterverhältnisse ist jedoch ohne die aktive Einbeziehung von Männern und eine grundlegende Reform des Bildungssystems nicht denkbar. Gleichstellung ist keine „Frauenfrage“, sondern eine gesamtgesellschaftliche Aufgabe.
	
	\subsection{Männer als aktive Verbündete (Male Allyship) im Alltag}
	
	Gleichstellung erreicht ihre volle Wirkung erst dann, wenn Männer diese nicht nur tolerieren, sondern aktiv als Verbündete (\textit{Male Allies}) vorantreiben. Wie eine Analyse der Rolle männlicher Verbündeter zeigt, beginnt Gleichstellung im privaten Alltag: Männer können durch die Übernahme von Haushaltspflichten, aktive Beteiligung an der Kindererziehung sowie durch die Unterstützung der Karriere ihrer Partnerinnen zu einer gerechteren Verteilung von Care-Arbeit beitragen \cite{male_allies}. Eine gerechte Verteilung von Sorge- und Erwerbsarbeit in Paar- und Familienbeziehungen trägt nachweislich nicht nur zur gesteigerten Lebensqualität bei, sondern reduziert auch das Risiko innerfamiliärer Gewalt, indem Machtasymmetrien innerhalb der Partnerschaft abgebaut werden \cite{bundesforum_gewalt}. So wird ein Klima der Gleichberechtigung begünstigt, das den Schutz vor Gewalt insbesondere gegen Frauen stärkt \cite{bundesforum_gewalt}.
	
	Am Arbeitsplatz können männliche Verbündete ihre weiblichen Kolleginnen unterstützen, indem sie sich für gleiche Bezahlung, faire Beförderungen und inklusive Richtlinien einsetzen \cite{male_allies}. Entscheidend ist auch die Intervention im sozialen Umfeld: Männer müssen ihre Stimme nutzen, um sexistische Bemerkungen, diskriminierendes Verhalten und unbewusste Vorurteile in Alltagssituationen zu kritisieren – egal ob in sozialen Kreisen, am Arbeitsplatz oder in Online-Plattformen \cite{male_allies}.
	
	Das Bundesforum Männer betont in seiner Positionierung zum Internationalen Tag gegen Gewalt an Frauen (25. November 2025), dass Gewaltschutz kein „Frauenthema“, sondern ein Gesellschaftsthema ist. Männer können und müssen dabei eine aktive Rolle einnehmen – als Verbündete und Mitgestaltende (Allyship) \cite{bundesforum_gewalt}. Dazu gehören: Hinschauen statt wegsehen und Solidarität mit Betroffenen zeigen, sexistische Sprache und abwertende Haltungen klar benennen sowie Zivilcourage zeigen \cite{bundesforum_gewalt}.
	
	Der HeForShe-Ansatz der Vereinten Nationen setzt genau hier an. Die Kampagne spricht Männer bewusst als Partner in der Gleichstellungsarbeit an \cite{bundesforum_gewalt}. Der HeForShe-Botschafter für UN Women Deutschland, Fikri Anıl Altıntaş, bringt es auf den Punkt: „Ich glaube fest daran, dass Männer Zugang zu feministischen Antworten in Zeiten der Unsicherheit erhalten, wenn wir mit ihnen sprechen und sie mitnehmen. Und dass der Austausch über Männlichkeit ein Weg dafür sein kann. Es geht nicht darum, sich selbst Feminist zu nennen, sondern um die Arbeit und Solidarität dahinter.“ \cite{heforshe_doku}.
	
	Praktische Formate wie der Workshop „Männer als Verbündete für Gleichstellung und Chancengleichheit“ der Landesfachstelle Männerarbeit Sachsen (2026) vermitteln konkrete Handlungsmöglichkeiten. Im Zentrum stehen die kritische Selbstreflexion, die Auseinandersetzung mit Privilegien und der Dialog auf Augenhöhe, um konkrete Ideen für solidarisches Handeln im Alltag, Beruf und gesellschaftlich-politischen Kontext zu entwickeln \cite{maenner_als_verbuendete_workshop}.
	
	\subsection{Bildungsreformen: Geschlechtergerechte Pädagogik als transformative Kraft}
	
	Das Bildungssystem ist ein zentraler Ort der Reproduktion, aber auch der Transformation von Geschlechterstereotypen. Die Stadt München hat mit ihrem im Dezember 2025 beschlossenen „Konzept für geschlechtergerechte Pädagogik für junge Menschen im Schulalter in Kindertageseinrichtungen, Schule und Ganztag“ eine wegweisende Strategie vorgelegt. Das Konzept verfolgt das übergeordnete Handlungsziel, Bildungsgerechtigkeit und Chancengleichheit für jeden jungen Menschen herzustellen, indem es Geschlechterstereotypen entgegenwirkt und geschlechtsspezifische Ungleichheiten abbaut \cite{muenchen_konzept_pm}. Im Zentrum stehen neu entwickelte Standards, mit denen die Qualität der geschlechtergerechten Pädagogik in Bildungseinrichtungen verbessert werden kann \cite{muenchen_konzept_pm}.
	
	Das Konzept adressiert mehrere Handlungsfelder: Organisationsentwicklung, Mitarbeitendenqualifikation, Innen- und Außenraumgestaltung sowie Ausstattung der Einrichtungen \cite{muenchen_konzept_pm}. Umfangreiche Checklisten erleichtern es den Pädagoginnen und Pädagogen, positive Alternativen zu bisherigen Praktiken zu finden \cite{muenchen_konzept_pm}. Es sensibilisiert die pädagogischen Fach- und Führungskräfte für stereotype Muster und unterstützt sie, das eigene Verhalten, die eigene Sprache und die Bildungsangebote zu reflektieren \cite{muenchen_konzept_artikel}. Bürgermeister Dominik Krause betont: „Eine vielfältige, offene Gesellschaft beginnt mit gleichen Chancen – für alle Kinder und Jugendlichen, unabhängig von Geschlecht, Herkunft oder sozialem Hintergrund.“ \cite{muenchen_konzept_artikel}.
	
	Auch die Förderung von Genderkompetenz in der Aus- und Fortbildung von Lehrkräften ist ein entscheidender Hebel. Die Bezirksregierung Arnsberg bietet Fortbildungen für Lehrkräfte aller Schulformen an, um geschlechtersensible Bildung als Zukunftsthema für die gesamte Gesellschaft zu multiplizieren \cite{arnsberg_fobi}. Die Angebote umfassen Grundlagen geschlechtersensibler Bildung, geschlechtliche und sexuelle Vielfalt, inklusive und diskriminierungsfreie Sprache, geschlechtersensibles Classroom-Management sowie geschlechtersensible berufliche Orientierung \cite{arnsberg_fobi}. Diese Maßnahmen unterstützen Lehrkräfte darin, Kinder und junge Erwachsene bestmöglich in ihrer Kompetenz- und Persönlichkeitsentwicklung zu fördern \cite{arnsberg_fobi}. Ein wichtiges Ziel ist es, die Berufswahl junger Menschen unvoreingenommen von Geschlechterklischees auf Basis individueller Begabungen zu ermöglichen, um so strukturellen Ungleichheiten wie Gender Pay Gap oder Gender Leadership Gap entgegenzuwirken \cite{muenchen_konzept_artikel}.
	
	\subsection{Fazit: Alltagshandeln und Bildung als komplementäre Strategien}
	
	Die aktive Beteiligung von Männern als Verbündete im Alltag und die Implementierung geschlechtergerechter Pädagogik im Bildungssystem sind zwei komplementäre Strategien, die einander wechselseitig verstärken: Wer in der Schule gelernt hat, Geschlechterstereotypen zu hinterfragen, wird als Erwachsener eher bereit sein, Verantwortung für Gleichstellung zu übernehmen. Bildung schafft die kognitive und emotionale Grundlage für gleichberechtigtes Handeln, während alltägliche Praktiken der Solidarität bestehende Strukturen von innen heraus verändern. Beide Ansätze adressieren die kulturellen und sozialen Ursachen von Ungleichheit – und sind damit unverzichtbare Ergänzungen zu rechtlichen und politischen Maßnahmen.
	
	\section{Weitere Herausforderungen und Schutzlücken}
	
	\subsection{Geschlechterungerechtigkeit im Gesundheitswesen}
	
	Die gesundheitliche Versorgung von Frauen in Deutschland weist erhebliche Lücken auf — von der Diagnostik von Endometriose und Wechseljahren über Zugangsbarrieren zu medizinisch sicheren Schwangerschaftsabbrüchen bis hin zu sicherer und selbstbestimmter Geburtshilfe. Die Deutsche Gesellschaft für Gynäkologie und Geburtshilfe (DGGG) sieht akuten Handlungsbedarf in mehreren Bereichen, darunter die unzureichende Berücksichtigung gynäkologischer Krebserkrankungen im Krankenhausstrukturgesetz, fehlende verbindliche Finanzierungsmodelle für die ärztliche Weiterbildung sowie Defizite in der geburtshilflichen Versorgung \cite{dggg_appell}.
	
	\subsection{Intersektionale Benachteiligungen}
	
	\textbf{Frauen mit Behinderungen} haben einen erschwerten Zugang zu medizinischer Versorgung und sind von häuslicher Gewalt besonders betroffen. Der Alternativbericht des Bündnisses Istanbul-Konvention weist darauf hin, dass vulnerable Gruppen — darunter Frauen mit Behinderungen, Geflüchtete, wohnungs- und obdachlose Frauen sowie trans-, inter- und queere Personen — durch institutionelle Hürden, diskriminierende Praktiken und fehlende Ressourcen kaum oder nur sehr erschwerten Zugang zum Hilfesystem haben \cite{bik_alternativbericht}.
	
	\textbf{Frauen mit Migrations- und Fluchthintergrund} sind mehrfach benachteiligt. Der Deutsche Frauenrat betont: Die Gleichstellung aller Frauen ist Voraussetzung für eine funktionierende Demokratie und gesellschaftlichen Zusammenhalt. Diskriminierung aufgrund der ethnischen Herkunft betrifft den gesamten Arbeitsprozess. Davon besonders betroffen sind Frauen, die ein Kopftuch tragen \cite{bmas_diskriminierung}.
	
	\textbf{LGBTQ+-Frauen} bleiben in vielen Bereichen strukturell benachteiligt, obwohl das Selbstbestimmungsgesetz einen wichtigen Fortschritt darstellt.
	
	\section{Der Europäische Fahrplan für Frauenrechte}
	
	Die Europäische Kommission hat am 7. März 2025 den \textit{Fahrplan für die Frauenrechte} angenommen. Die Mitteilung (COM(2025) 97 final) enthält eine langfristige Vision für die Gleichstellung der Geschlechter, die auf den zentralen Grundsätzen und politischen Zielen im Anhang der Erklärung der Grundsätze für eine geschlechtergerechte Gesellschaft beruht \cite{eu_fahrplan}.
	
	In der Mitteilung wird hervorgehoben, dass "die Gleichheit von Frauen und Männern ein Grundrecht ist und als Grundwert seit 1957 mit den Römischen Verträgen im EU-Recht verankert" sei. Die Kommission führt aus, dass die Gleichstellung der Geschlechter für Wettbewerbsfähigkeit, Demokratie und soziale Gerechtigkeit unabdingbar ist \cite{eu_fahrplan_bremen}. Prognosen zufolge könnte eine Verbesserung bei der Gleichstellung der Geschlechter das Wachstum des Pro-Kopf-BIP der EU bis 2050 um 6,1\% bis 9,6\% heben \cite{eu_fahrplan_bremen}.
	
	\section{Schlussfolgerungen und Handlungsempfehlungen}
	
	Auf Basis der Analyse ergeben sich folgende politische Handlungsempfehlungen:
	
	\begin{enumerate}[label=\arabic*.]
		\item \textbf{Konsequente Umsetzung der Istanbul-Konvention:} Verpflichtende, bundesweit einheitliche und ausreichend finanzierte Schutzstrukturen müssen geschaffen werden — mit besonderem Fokus auf vulnerable Gruppen. Der Alternativbericht des BIK zeigt, dass ein grundlegender Paradigmenwechsel von Leuchtturm-, Pilot- und Modellprojekten hin zu flächendeckenden, nachhaltig finanzierten und diskriminierungsfreien Strukturen im Gewaltschutz erforderlich ist \cite{bik_alternativbericht}.
		
		\item \textbf{Verbindliche Regelungen für Parität in Führungspositionen:} Die anhaltende Stagnation beim Frauenanteil in Führungspositionen erfordert gesetzlich verankerte Regelungen für Parität in den Führungspositionen von Wirtschaft, Verwaltung, Wissenschaft und Kultur.
		
		\item \textbf{Bekämpfung digitaler Gewalt:} Der DSA muss konsequent umgesetzt und weiterentwickelt werden. Digitale Gewalt ist als Form geschlechtsspezifischer Gewalt rechtlich anzuerkennen und mit wirksamen Sanktionen zu belegen. Die EU-Gewaltschutz-Richtlinie 2024 bietet hierfür einen Rahmen, der nun zügig in nationales Recht umgesetzt werden muss.
		
		\item \textbf{Stärkung der politischen Partizipation:} Paritätisch besetzte Listenmandate bei Wahlen und die paritätische Vertretung in allen Entscheidungsorganen sind notwendig, um den rückläufigen Frauenanteil im Parlament umzukehren.
		
		\item \textbf{Geschlechtergerechte Gesundheitsversorgung:} Eine flächendeckende, barrierefreie und geschlechtergerechte Gesundheitsversorgung muss sichergestellt werden, wie dies von der Deutschen Gesellschaft für Gynäkologie und Geburtshilfe gefordert wird \cite{dggg_appell}.
		
		\item \textbf{Intersektionale Perspektiven institutionalisieren:} Gleichstellungspolitik muss die Mehrfachbenachteiligungen von Frauen mit Behinderungen, Migrationshintergrund und LGBTQ+-Frauen systematisch adressieren.
		
		\item \textbf{Feministische KI-Strategie entwickeln:} Auf Grundlage des EU AI Act müssen verbindliche Geschlechterverträglichkeitsprüfungen für KI-Systeme eingeführt, der Gender Data Gap geschlossen und die Beteiligung von Frauenverbänden in KI-Ethikgremien sichergestellt werden \cite{eu_ai_act_frauenrat}.
		
		\item \textbf{Männer als Verbündete aktivieren:} Die gezielte Ansprache von Männern als aktive Partner in der Gleichstellungsarbeit (Male Allyship) ist systematisch auszubauen. Dazu gehören öffentlichkeitswirksame Kampagnen (wie HeForShe), niedrigschwellige Workshop-Angebote zur kritischen Selbstreflexion sowie die Integration von Gleichstellungsthemen in die männliche Berufs- und Gesundheitsberatung.
		
		\item \textbf{Geschlechtergerechte Bildung verankern:} Das Münchner Konzept für geschlechtergerechte Pädagogik sollte als Best-Practice-Beispiel bundesweit adaptiert werden. Geschlechtergerechtigkeit ist als Querschnittsaufgabe in allen Bildungsplänen zu verankern. Genderkompetenz muss verpflichtender Bestandteil der Lehrkräfteaus- und -fortbildung werden – von der frühkindlichen Bildung bis zur Hochschule.
	\end{enumerate}
	
	\section{Fazit}
	
	Deutschland hat bei der Gleichstellung der Geschlechter in zentralen Bereichen erheblichen Nachholbedarf — beim Gewaltschutz, auf dem Arbeitsmarkt, in Führungsetagen, in der Politik, im digitalen Raum und zunehmend auch im Bereich der Künstlichen Intelligenz. Die Demokratie bleibt unvollständig, solange die Mehrheit der Bevölkerung strukturell benachteiligt wird. Die Europäische Union hat mit ihrem Fahrplan für die Frauenrechte (März 2025) eine langfristige Vision vorgelegt, die insbesondere auch neue Herausforderungen wie geschlechtsspezifische Verzerrungen, Diskriminierung und Gewalt im Internet adressiert. Die Umsetzung dieser Vision erfordert nun politischen Willen, ausreichende Ressourcen und eine konsequente feministische Ausrichtung aller Politikfelder.
	
	\newpage
	
	\begin{thebibliography}{99}
		
		\bibitem{gewalthilfegesetz}
		Nomos: \emph{Gewalthilfegesetz (GewHG)}. 2025.
		URL: \url{https://www.nomos.de/wp-content/uploads/2025/03/Beitrag-Gewalthilfegesetz.pdf}
		
		\bibitem{bik_alternativbericht}
		Frauenhauskoordinierung e.V.: \emph{Alternativbericht zur Istanbul-Konvention 2025: Deutschland wird seinen Verpflichtungen nicht gerecht}. 2025.
		URL: \url{https://www.frauenhauskoordinierung.de/publikationen/detail/alternativbericht-istanbul-konvention-2025}
		
		\bibitem{gpg_2025}
		Tagesschau: \emph{Gender Pay Gap 2025: Lohnlücke bleibt bei 16 Prozent}. 2026.
		URL: \url{https://www.tagesschau.de/wirtschaft/arbeitsmarkt/gender-pay-gap-lohnluecke-frauen-100.html}
		
		\bibitem{gpg_bereinigt}
		Deutschlandfunk: \emph{Für die gleiche Arbeit 6 \% weniger Lohn – Was hilft gegen die schlechtere Bezahlung von Frauen?}. 2026.
		URL: \url{https://www.deutschlandfunk.de/gender-pay-gap-equal-pay-day-entgeltgleichheit-100.html}
		
		\bibitem{gpg_ost_west}
		Tagesschau: \emph{Frauen verdienen weiterhin weniger als Männer}. 2025.
		URL: \url{https://www.tagesschau.de/wirtschaft/arbeitsmarkt/gender-pay-gap-lohnluecke-frauen-100.html}
		
		\bibitem{frauenanteil_bundestag}
		Bundesstiftung Gleichstellung: \emph{Rückschritt für die Gleichstellung: Warum der Frauenanteil im Bundestag sinkt}. 2025.
		URL: \url{https://www.bundesstiftung-gleichstellung.de/2025/03/31/rueckschritt-fuer-die-gleichstellung-warum-der-frauenanteil-im-bundestag-sinkt/}
		
		\bibitem{dsa_luecken}
		Deutscher Juristinnenbund e.V.: \emph{DSA gilt seit 17.02.24, Schutz vor digitaler Gewalt im Netz weiterhin lückenhaft}. 2024.
		URL: \url{https://www.djb.de/fileadmin/user_upload/presse/pressemitteilungen/pm24-15_DSA.pdf}
		
		\bibitem{eu_gewalt_richtlinie}
		Europäische Union: \emph{Verhütung und Bekämpfung von Gewalt gegen Frauen und häuslicher Gewalt in der Europäischen Union}. 2024.
		URL: \url{https://eur-lex.europa.eu/DE/legal-content/summary/preventing-and-combating-violence-against-women-and-domestic-violence-throughout-the-european-union.html}
		
		\bibitem{eu_gewalt_umsetzung}
		Europäische Union: \emph{Richtlinie (EU) 2024/1385 zur Bekämpfung von Gewalt gegen Frauen und häuslicher Gewalt}. 2024.
		URL: \url{https://eur-lex.europa.eu/legal-content/DE/TXT/?uri=CELEX\%3A32024L1385}
		
		\bibitem{eu_fahrplan}
		Europäische Kommission: \emph{Mitteilung der Kommission an das Europäische Parlament, den Rat, den Europäischen Wirtschafts- und Sozialausschuss und den Ausschuss der Regionen: Fahrplan für die Frauenrechte (COM(2025) 97 final)}. 2025.
		URL: \url{https://eur-lex.europa.eu/legal-content/DE/TXT/?uri=COM\%3A2025\%3A97\%3AFIN}
		
		\bibitem{eu_fahrplan_bremen}
		Senatskanzlei Bremen: \emph{Ein Fahrplan für Frauenrechte}. 2025.
		URL: \url{https://sd.bremische-buergerschaft.de/sdnetrim/UGhVM0hpd2NXNFdFcExjZbsu-KIYVhtlzPacKXzfzcOG6Ed2mA5dRBBal2-5wW8B/Beschlussvorlage_Ausschuesse-Deputationen_VL_21-5061.pdf}
		
		\bibitem{dggg_appell}
		Deutsche Gesellschaft für Gynäkologie und Geburtshilfe e.V.: \emph{Reform-Appell an die neue Bundesregierung für eine bessere Gesundheitsversorgung von Frauen in Deutschland}. 2025.
		URL: \url{https://www.dggg.de/presse/pressemitteilungen-und-nachrichten/reform-appell-an-die-neue-bundesregierung-fuer-eine-bessere-gesundheitsversorgung-von-frauen-in-deutschland}
		
		\bibitem{bmas_diskriminierung}
		Bundesministerium für Arbeit und Soziales: \emph{Teilhabebefragung von Migrantinnen, Migranten und Schutzsuchenden}. 2025.
		URL: \url{https://www.bmas.de/SharedDocs/Downloads/DE/Publikationen/Forschungsberichte/fb-672-migrant-innen-schutzsuchende-teilhabebefragung.pdf}
		
		\bibitem{thws_gehaltsbias}
		Technische Hochschule Würzburg-Schweinfurt: \emph{Künstliche Intelligenz gibt Frauen schlechtere Gehaltsratschläge – THWS-Studie deckt Voreingenommenheit bei Sprachmodellen wie ChatGPT auf}. 2025.
		URL: \url{https://idw-online.de/de/news855020}
		
		\bibitem{thws_gehaltsbias_pm}
		idw - Informationsdienst Wissenschaft: \emph{Künstliche Intelligenz gibt Frauen schlechtere Gehaltsratschläge (Pressemitteilung)}. 2025.
		URL: \url{https://idw-online.de/de/news855020}
		
		\bibitem{bka_deepfakes}
		Bundeskriminalamt (BKA): \emph{Von politischer Desinformation zu Nudification Apps – Deepfakes und digitale Gewalt im regulatorischen Kontext}. Präsentation, Forum KI 2025, Wiesbaden, 13. Mai 2025.
		URL: \url{https://www.bka.de/SharedDocs/Downloads/DE/Publikationen/ForumKI/ForumKI2025/Präsentation/10_Präsentation_Benning.pdf}
		
		\bibitem{messer_deepfake}
		Messer, Lilli: \emph{Visuelle Verfügbarkeit über Deepfake – eine digitale Form sexualisierter Gewalt}. Universität Duisburg-Essen, 2025.
		URL: \url{https://duepublico2.uni-due.de/servlets/MCRFileNodeServlet/duepublico_derivate_00083327/Messer_deepfake_sexualisierte_gewalt.pdf}
		
		\bibitem{eu_ai_act_frauenrat}
		Deutscher Frauenrat: \emph{Für eine geschlechtergerechte Künstliche Intelligenz: Entwicklung einer feministischen KI-Strategie auf der Grundlage des „EU Artificial Intelligence Act“ (EU AI Act)}. Beschluss vom 27. Juni 2025.
		URL: \url{https://www.frauenrat.de/fuer-eine-geschlechtergerechte-kuenstliche-intelligenz-entwicklung-einer-feministischen-ki-strategie-auf-der-grundlage-des-eu-artificial-intelligence-act-eu-ai-act/}
		
		\bibitem{europarat_empfehlung_ki}
		Europarat: \emph{Empfehlung zur Gleichstellung und künstlichen Intelligenz}. 2025.
		URL: \url{https://eur-lex.europa.eu/legal-content/DE/TXT/?uri=CELEX\%3A32025H0350}
		
		\bibitem{male_allies}
		Daily Tribune: \emph{Male allies: the role of men in supporting Gender Equality}. 8. März 2025.
		URL: \url{https://tribune.net.ph/2025/03/08/male-allies-the-role-of-men-in-supporting-gender-equality}
		
		\bibitem{bundesforum_gewalt}
		Bundesforum Männer: \emph{Internationaler Tag gegen Gewalt an Frauen: Gewalt geht uns alle an}. 25. November 2025.
		URL: \url{https://bundesforum-maenner.de/position/internationaler-tag-gegen-gewalt-an-frauen-gewalt-geht-uns-alle-an/}
		
		\bibitem{heforshe_doku}
		HeForShe / UN Women: \emph{HeForShe Advocate for UN Women Germany Champions Gender Equality in ARTE's New Documentary}. 2025.
		URL: \url{https://www.heforshe.org/fr/node/511}
		
		\bibitem{maenner_als_verbuendete_workshop}
		Landesfachstelle Männerarbeit Sachsen: \emph{Praxis-Workshop „Männer als Verbündete für Gleichstellung und Chancengleichheit – Male Allyship“}. 2026.
		URL: \url{https://www.fachstelle-maennerarbeit.de/event/male-allyship-praxisworkshop/}
		
		\bibitem{muenchen_konzept_pm}
		Landeshauptstadt München: \emph{Konzept für geschlechtergerechte Pädagogik für junge Menschen im Schulalter in Kindertageseinrichtungen, Schule und Ganztag}. Dezember 2025.
		URL: \url{https://risi.muenchen.de/risi/dokument/v/9412961}
		
		\bibitem{muenchen_konzept_artikel}
		Landeshauptstadt München, Rathaus Umschau: \emph{Stadt präsentiert Konzept zu geschlechtergerechter Pädagogik}. 4. Dezember 2025.
		URL: \url{https://ru.muenchen.de/2025/232/Stadt-praesentiert-Konzept-zu-geschlechtergerechter-Paedagogik-121786}
		
		\bibitem{arnsberg_fobi}
		Bezirksregierung Arnsberg: \emph{Genderkompetenz stärken: Fortbildungen für eine geschlechtersensible Bildung}. Schuljahr 2024/25.
		URL: \url{https://www.bra.nrw.de/bildung-schule/lehrkraeftebildung/fortbildung/fortbildung-fuer-lehrkraefte-berufskollegs/aktuelles-zur-fortbildung-fuer-lehrkraefte-berufskollegs/genderkompetenz-staerken-fortbildungen-fuer-eine-geschlechtersensible-bildung}
		
	\end{thebibliography}
	
\end{document}